\section{More functoriality of presheaves}
\label{section-more-functoriality-PSh}

\noindent
In this section we revisit the material of
Section \ref{section-functoriality-PSh}.
Let $u : \mathcal{C} \to \mathcal{D}$ be a functor between categories.
Recall that
$$
u^p :
\textit{PSh}(\mathcal{D})
\longrightarrow
\textit{PSh}(\mathcal{C})
$$
is the functor that associates to $\mathcal{G}$ on $\mathcal{D}$ the presheaf
$u^p\mathcal{G} = \mathcal{G} \circ u$. It turns out that this functor
not only has a left adjoint (namely $u_p$) but also a right adjoint.

\medskip\noindent
Namely, for any $V \in \Ob(\mathcal{D})$
we define a category ${}_V\mathcal{I} = {}_V^u\mathcal{I}$.
Its objects are pairs $(U, \psi : u(U) \to V)$.
Note that the arrow is in the opposite direction from the arrow
we used in defining the category $\mathcal{I}_V^u$ in
Section \ref{section-functoriality-PSh}.
A morphism $(U, \psi) \to (U', \psi')$ is given
by a morphism $\alpha : U \to U'$ such that
$\psi = \psi' \circ u(\alpha)$.
In addition, given any presheaf of sets $\mathcal{F}$
on $\mathcal{C}$ we introduce the functor
${}_V\mathcal{F} : {}_V\mathcal{I}^{opp} \to \textit{Sets}$,
which is defined by the rule
${}_V\mathcal{F}(U, \psi) = \mathcal{F}(U)$. We define
$$
{}_pu(\mathcal{F})(V) := \lim_{{}_V\mathcal{I}^{opp}} {}_V\mathcal{F}
$$
As a limit there are projection maps
$c(\psi) : {}_pu(\mathcal{F})(V) \to \mathcal{F}(U)$
for every object $(U, \psi)$ of ${}_V\mathcal{I}$.
In fact,
$$
{}_pu(\mathcal{F})(V)
=
\left\{
\begin{matrix}
\text{collections }
s_{(U, \psi)} \in \mathcal{F}(U) \\
\forall \beta : (U_1, \psi_1) \to (U_2, \psi_2)
\text{ in }{}_V\mathcal{I} \\
\text{ we have } \beta^*s_{(U_2, \psi_2)} = s_{(U_1, \psi_1)}
\end{matrix}
\right\}
$$
where the correspondence is given by $s \mapsto s_{(U, \psi)} = c(\psi)(s)$.
We leave it to the reader to define the restriction mappings
${}_pu(\mathcal{F})(V) \to {}_pu(\mathcal{F})(V')$
associated to any morphism $V' \to V$ of $\mathcal{D}$.
The resulting presheaf will be denoted ${}_pu\mathcal{F}$.

\begin{lemma}
\label{lemma-recover-pu}
There is a canonical map
${}_pu\mathcal{F}(u(U)) \to \mathcal{F}(U)$,
which is compatible with restriction maps.
\end{lemma}

\begin{proof}
This is just the projection map $c(\text{id}_{u(U)})$ above.
\end{proof}

\noindent
Note that any map of presheaves $\mathcal{F} \to \mathcal{F}'$
gives rise to compatible systems of maps between functors
${}_V\mathcal{F} \to {}_V\mathcal{F}'$, and hence to a map
of presheaves ${}_pu\mathcal{F} \to {}_pu\mathcal{F}'$. In other
words, we have defined a functor
$$
{}_pu :
\textit{PSh}(\mathcal{C})
\longrightarrow
\textit{PSh}(\mathcal{D})
$$

\begin{lemma}
\label{lemma-adjoints-pu}
The functor ${}_pu$ is a right adjoint to the functor $u^p$.
In other words the formula
$$
\Mor_{\textit{PSh}(\mathcal{C})}(u^p\mathcal{G}, \mathcal{F})
=
\Mor_{\textit{PSh}(\mathcal{D})}(\mathcal{G}, {}_pu\mathcal{F})
$$
holds bifunctorially in $\mathcal{F}$ and $\mathcal{G}$.
\end{lemma}

\begin{proof}
This is proved in exactly the same way as the proof
of Lemma \ref{lemma-adjoints-u}. We note that the map
$u^p{}_pu \mathcal{F} \to \mathcal{F}$ from
Lemma \ref{lemma-recover-pu} is the map that
is used to go from the right to the left.

\medskip\noindent
Alternately, think
of a presheaf of sets $\mathcal{F}$ on $\mathcal{C}$ as a presheaf
$\mathcal{F}'$ on $\mathcal{C}^{opp}$ with values in $\textit{Sets}^{opp}$,
and similarly on $\mathcal{D}$.
Check that $({}_pu \mathcal{F})' = u_p(\mathcal{F}')$,
and that $(u^p\mathcal{G})' = u^p(\mathcal{G}')$.
By Remark \ref{remark-functoriality-presheaves-values}
we have the adjointness of $u_p$ and $u^p$ for
presheaves with values in $\textit{Sets}^{opp}$.
The result then follows formally from this.
\end{proof}

\noindent
Thus given a functor $u : \mathcal{C} \to \mathcal{D}$ of categories
we obtain a sequence of functors
$$
u_p, u^p, {}_pu
$$
between categories of presheaves where in each consecutive pair the
first is left adjoint to the second.

\begin{lemma}
\label{lemma-adjoint-functors}
Let $u : \mathcal{C} \to \mathcal{D}$ and $v : \mathcal{D} \to \mathcal{C}$
be functors of categories. Assume that $v$ is right adjoint to $u$.
Then we have
\begin{enumerate}
\item $u^ph_V = h_{v(V)}$ for any $V$ in $\mathcal{D}$,
\item the category $\mathcal{I}^v_U$ has an initial object,
\item the category ${}_V^u\mathcal{I}$ has a final object,
\item ${}_pu = v^p$, and
\item $u^p = v_p$.
\end{enumerate}
\end{lemma}

\begin{proof}
Proof of (1). Let $V$ be an object of $\mathcal{D}$. We have
$u^ph_V = h_{v(V)}$ because
$u^ph_V(U) = \Mor_\mathcal{D}(u(U), V) = \Mor_\mathcal{C}(U, v(V))$
by assumption.

\medskip\noindent
Proof of (2). Let $U$ be an object of $\mathcal{C}$. Let
$\eta : U \to v(u(U))$ be the map adjoint to the map
$\text{id} : u(U) \to u(U)$. Then we claim $(u(U), \eta)$
is an initial object of $\mathcal{I}_U^v$. Namely, given
an object $(V, \phi : U \to v(V))$ of $\mathcal{I}_U^v$
the morphism $\phi$ is adjoint to a map $\psi : u(U) \to V$
which then defines a morphism $(u(U), \eta) \to (V, \phi)$.

\medskip\noindent
Proof of (3). Let $V$ be an object of $\mathcal{D}$. Let
$\xi : u(v(V)) \to V$ be the map adjoint to the map
$\text{id} : v(V) \to v(V)$. Then we claim $(v(V), \xi)$
is a final object of ${}_V^u\mathcal{I}$. Namely, given
an object $(U, \psi : u(U) \to V)$ of ${}_V^u\mathcal{I}$
the morphism $\psi$ is adjoint to a map $\phi : U \to v(V)$
which then defines a morphism $(U, \psi) \to (v(V), \xi)$.

\medskip\noindent
Hence for any presheaf $\mathcal{F}$ on $\mathcal{C}$ we have
\begin{eqnarray*}
v^p\mathcal{F}(V)
& = &
\mathcal{F}(v(V)) \\
& = &
\Mor_{\textit{PSh}(\mathcal{C})}(h_{v(V)}, \mathcal{F}) \\
& = &
\Mor_{\textit{PSh}(\mathcal{C})}(u^ph_V, \mathcal{F}) \\
& = &
\Mor_{\textit{PSh}(\mathcal{D})}(h_V, {}_pu\mathcal{F}) \\
& = &
{}_pu\mathcal{F}(V)
\end{eqnarray*}
which proves part (4). Part (5) follows by the uniqueness of adjoint functors.
\end{proof}




